\section{Introduction}\label{sec:introduction}
Neural approximations to dynamic economic decisions require a criterion that is linked to the value of the policy actually delivered. Average equation residuals need not provide that link. A small error at frequently visited training states can coexist with a large error near a stopping face, and improving a value witness need not improve the policy with which it is paired. These distinctions become operational when an algorithm must return a complete error bound rather than a sampled diagnostic.

This paper develops certificate-aligned neural Bellman operators. The primary algorithm minimizes a differentiable envelope of two signed Bellman residuals on a complete state--time cover. An independent directed-arithmetic verifier checks the proposed policy and witness against the original economic model. Every prescribed proposal is retained; an incumbent is replaced only when its complete certificate decreases. A quantitative theorem connects the trainable envelope to the rigorous policy-loss bound, separating the smoothing allowance from a checkable numerical correction. The theorem concerns the loss actually used in the computation, rather than an unexecuted policy-iteration idealization.

Stopping geometry is part of this construction. In our consumption--investment economy with costly preference adjustment, continuation value has a strict jump exceeding $6.37$ utility units at an inward upper-wealth face. A critic forced to equal immediate settlement on every face therefore cannot acquire a vanishing uniform positive residual. A signed accessibility-weighted comparison permits the critic to dominate settlement at a face that the candidate cannot reach, while retaining every originally admissible action in the optimal upper comparison. This corrects an approximation-space inconsistency without modifying preferences, controls, or the stopping contract.

The original-economy experiment compares two generations of training on that unchanged problem. Under the historical sampled objective, extending training from 800 to 2,400 updates worsens all ten certificates, even though sampled residual mean squared error falls for nine seeds. Under the complete-cover objective, all ten accessibility-aware trajectories have decreasing raw certificates at 0, 100, and 400 updates. Their final bounds range from $7.27832$ to $7.34367$, compared with initial bounds from $26.3005$ to $26.6668$. A paired ten-seed all-face experiment, unrestricted state-space comparators with method-specific witnesses, and a fixed-network refinement study identify what causes this improvement and what it does not establish. In particular, the $.01$ original-economy target is unchanged and is not attained. A certified fixed-witness floor and actor--critic interchange show why neither further arithmetic precision nor a monotone incumbent alone establishes accurate policy learning.

A second, nonlinear application supplies an explicit accuracy--work result on a genuinely vector-valued state. A heterogeneous multi-product inventory planner has 12 decision dates, coupled nonquadratic costs, and up to 128 state coordinates. The neural map takes the full initial state and returns a control plan, without optimal-value or optimal-control labels. Strong convexity yields a state-uniform certificate for exact-real gradient corrections to that plan. At 128 coordinates, the total-cost bound falls from $0.124820$ after 16 corrections to $0.000359609$ after 24 and below $1.037\times10^{-6}$ after 32. This is an unknown-value nonlinear planning problem, not the two-mode Riccati example retained for validation. Zero-start gradient, accelerated-gradient, and quasi-Newton alternatives are evaluated with the same independent cost-loss certificate. They remain strong competitors; the experiment establishes an accurate neural-initialized solver, not a demonstrated neural speed advantage.

Economic interpretation also requires keeping numerical objects distinct. The sharp deterministic time-control/dual library in the original economy is not a neural result. Its price-uniform regret below $0.009982$ can be offset by a rigorously checked, budget-financed initial-wealth increment of $.0078$ at initial wealth $1.25$, using an explicitly adjusted feasible consumption policy. This sufficient compensation measure respects the original consumption cap and stopping contract. The earlier externally financed top-up remains a secondary normalization. Neither measure converts the much larger original neural certificates into sharp welfare claims.

\subsection{Contribution and relation to the literature}
Stopped verification, Bellman contraction, strong-convexity error bounds, and gradient contraction are classical ingredients. \citet{Judd1998} and \citet{KushnerDupuis2001} develop numerical dynamic programming and controlled Markov-chain approximation. \citet{BarlesSouganidis1991} establishes a fundamental monotone-approximation framework, and \citet{BrummScheidegger2017} develops adaptive sparse-grid methods for computational economics. Information and control relaxations provide dual comparisons \citep{BrownSmithSun2010,BrownSmith2011}. Neural dynamic-control and differential-equation methods include \citet{HanE2016,HanJentzenE2018}; the historical external comparison uses \citet{CaiFangZhou2025}.

Our contribution is the combination of an identified stopping-trace inconsistency, an executed complete-cover neural objective with a checkable policy-loss implication, and a nonlinear control specialization with an explicit state-uniform work--accuracy bound. The new objective theorem does not claim global convergence of Adam, and the nonlinear specialization does not claim to remove the difficulty of arbitrary stochastic high-dimensional Bellman equations. The original-economy upper comparison always retains the full control opportunity. The arithmetic implementation uses MPFR \citep{FousseEtAl2007}; shared analytic formulas, stopping arguments, and program logic remain part of its trusted mathematical base.

The paper presents the model and verification foundations first, followed by the stopping correction, the certificate-aligned algorithm, the original-economy evidence, and the nonlinear control application. Historical coefficient transport, state--price continuation, external holdouts, adverse training outcomes, and specialized dual calculations remain available in the main text and supplement. All new numerical statements are tied to immutable source and result objects; all 60 prescribed original-economy checkpoints are included in the new objective audit.
